\documentclass[aps,prd,reprint,nofootinbib,superscriptaddress]{revtex4-2}
\usepackage{amsmath,amssymb,mathtools,bm}
\usepackage{microtype}
\usepackage[hidelinks]{hyperref}
\usepackage[T1]{fontenc}
\usepackage{newtxtext,newtxmath}

\newcommand{\SO}{\operatorname{SO}}
\newcommand{\Hol}{\operatorname{Hol}}
\newcommand{\R}{\mathbb R}

\begin{document}

\title{Compatibility Without Flatness:\\Exact Gluing with Nonzero Curvature, Holonomy, and Proper-Time Difference}

\author{A. V. Coppa}
\affiliation{Coppalini Architecture}
\date{2 September 2026}

\begin{abstract}
We give an explicit Lorentzian $\SO(2)$ gauge-field witness in which local bundle descriptions glue exactly while physical curvature, global holonomy, and proper-time difference remain independently readable. A nonnegative gluing functional is built from section, connection, and transition-map defects; its zero locus is exactly the simultaneous gluing condition. On $\R_t\times S^1\times\R_x\times\R_y$, the connection $A_a=a\eta+f(t-x/c)\,dy$ has nonzero source-free Maxwell curvature independent of $a$, while its holonomy around the circle is freely tunable by $a$. Two timelike paths with common endpoints have unequal proper time independently of both. The construction isolates one point: exact closure of local descriptions is a compatibility statement and need not be a flatness condition on the physical connection.
\end{abstract}

\maketitle

\section{Local bundle data and the compatibility functional}

Let $(M,g)$ be a time-oriented Lorentzian four-manifold and let $h$ be an auxiliary Riemannian metric used only to form positive norms. Let $E\to M$ be an oriented rank-two real vector bundle with structure group $\SO(2)$ and positive fiber metric $h_E$. On a finite trivializing cover $\{U_\alpha\}$, choose local section representatives $q_\alpha$, transition maps
\begin{equation}
U_{\alpha\beta}=e^{\chi_{\alpha\beta}C_1}\in\SO(2),
\qquad
C_1=
\begin{pmatrix}
0&-1\\
1&0
\end{pmatrix},
\label{eq:transitions}
\end{equation}
and local scalar connection one-forms $A_\alpha\in\Omega^1(U_\alpha)$. The standard connection transformation law is
\begin{equation}
A_\beta=A_\alpha-d\chi_{\alpha\beta}.
\label{eq:connlaw}
\end{equation}
Define the local compatibility defects
\begin{align}
\Delta^q_{\alpha\beta}
&:=q_\beta-U_{\alpha\beta}q_\alpha,
\\
\Delta^A_{\alpha\beta}
&:=A_\beta-A_\alpha+d\chi_{\alpha\beta},
\\
\Delta^U_{\alpha\beta\gamma}
&:=U_{\alpha\beta}U_{\beta\gamma}U_{\gamma\alpha}-I.
\label{eq:defects}
\end{align}
Assume continuity and square-integrability on the relevant overlaps. For fixed $\kappa_A,\kappa_U>0$, set
\begin{align}
\mathcal F_q
&:=\sum_{\alpha<\beta}\int_{U_{\alpha\beta}}
\|\Delta^q_{\alpha\beta}\|_{h_E}^2\,d\mathrm{vol}_h,
\\
\mathcal F_A
&:=\sum_{\alpha<\beta}\int_{U_{\alpha\beta}}
\|\Delta^A_{\alpha\beta}\|_{h}^2\,d\mathrm{vol}_h,
\\
\mathcal F_U
&:=\sum_{\alpha<\beta<\gamma}\int_{U_{\alpha\beta\gamma}}
\|\Delta^U_{\alpha\beta\gamma}\|^2\,d\mathrm{vol}_h,
\\
\mathcal F_{\Lambda}
&:=\mathcal F_q+\kappa_A\mathcal F_A+\kappa_U\mathcal F_U.
\label{eq:functional}
\end{align}

\paragraph{Proposition 1 (exact gluing).}
Under the stated regularity assumptions,
\begin{equation}
\boxed{
\mathcal F_{\Lambda}=0
\iff
\Delta^q_{\alpha\beta}=0,
\quad
\Delta^A_{\alpha\beta}=0,
\quad
\Delta^U_{\alpha\beta\gamma}=0
}
\label{eq:zero}
\end{equation}
on every relevant overlap.

\emph{Proof.}
Each term in Eq.~\eqref{eq:functional} is nonnegative. Vanishing defects give $\mathcal F_{\Lambda}=0$. Conversely, a zero finite sum with positive coefficients forces every integral to vanish. Continuity then forces each squared norm to vanish pointwise on its overlap. \hfill$\square$

The auxiliary metric affects numerical values away from zero but does not alter the common zero locus used as the gluing criterion.

\section{Explicit Lorentzian witness}

Take
\begin{equation}
M=\R_t\times S^1_\vartheta\times\R_x\times\R_y
\label{eq:manifold}
\end{equation}
and let $\eta$ be the global closed one-form on the circle normalized by
\begin{equation}
d\eta=0,
\qquad
\oint_{S^1}\eta=2\pi.
\end{equation}
For constants $c,\rho>0$, define
\begin{equation}
g=-c^2dt^2+\rho^2\eta^2+dx^2+dy^2.
\label{eq:metric}
\end{equation}
Use the trivial bundle $E=M\times\R^2$. Let $f\in C^2(\R)$ with $f'\not\equiv0$, set
\begin{equation}
u=t-\frac{x}{c},
\end{equation}
and define the scalar connection family
\begin{equation}
\boxed{A_a=a\eta+f(u)\,dy,\qquad a\in\R.}
\label{eq:connection}
\end{equation}
The corresponding matrix connection is $\widehat A_a=A_aC_1$.

Its curvature is
\begin{equation}
\boxed{F=dA_a=f'(u)\,du\wedge dy,}
\label{eq:curvature}
\end{equation}
which is independent of $a$ and nonzero whenever $f'\neq0$. The nonzero coordinate components are
\begin{equation}
F_{ty}=f'(u),
\qquad
F_{xy}=-\frac1c f'(u).
\end{equation}
Because the coefficients of Eq.~\eqref{eq:metric} are constant,
\begin{equation}
\nabla_\mu F^{\mu y}
=-\frac1{c^2}f''(u)+\frac1{c^2}f''(u)=0,
\end{equation}
and the other divergences vanish. Thus
\begin{equation}
dF=0,
\qquad
d\star_gF=0,
\qquad
F\neq0,
\label{eq:maxwell}
\end{equation}
so the witness is a nonzero source-free Maxwell field.

Because Eq.~\eqref{eq:connection} is globally defined, restrict it and any global section to a finite trivializing cover using identity transition maps. Every defect in Eq.~\eqref{eq:defects} vanishes, and therefore
\begin{equation}
\boxed{\mathcal F_{\Lambda}=0}
\label{eq:glued}
\end{equation}
while Eq.~\eqref{eq:curvature} remains nonzero.

\section{Holonomy and proper time remain independent}

For a loop $\ell$ winding once around $S^1$ with $dy=0$, Abelian parallel transport gives
\begin{equation}
\boxed{
\Hol_{A_a}(\ell)
=
\exp(-2\pi a C_1).
}
\label{eq:holonomy}
\end{equation}
Varying $a$ spans $\SO(2)$ while the curvature Eq.~\eqref{eq:curvature} remains fixed. Local curvature and global loop holonomy are therefore independently tunable in this family.

Now fix
\begin{equation}
p=(0,[0],0,0),
\qquad
q=(T,[0],0,0),
\qquad
T>\frac{2\pi\rho}{c},
\end{equation}
and define
\begin{align}
\gamma_0(t)&=(t,[0],0,0),\\
\gamma_1(t)&=\left(t,\left[\frac{2\pi t}{T}\right],0,0\right),
\qquad 0\le t\le T.
\end{align}
Both paths are timelike and have the same endpoints. Their proper times are
\begin{equation}
\tau[\gamma_0]=T,
\qquad
\tau[\gamma_1]
=T\sqrt{1-\left(\frac{2\pi\rho}{cT}\right)^2},
\end{equation}
so
\begin{equation}
\boxed{\Delta\tau:=\tau[\gamma_0]-\tau[\gamma_1]>0.}
\label{eq:proper}
\end{equation}
The proper-time difference depends on $(\rho,T)$, the holonomy on $a$, and the local Maxwell profile on $f$. These data can be varied independently within the explicit construction.

\section{Compatibility theorem and boundary}

Combining Eqs.~\eqref{eq:curvature}, \eqref{eq:glued}, \eqref{eq:holonomy}, and \eqref{eq:proper} gives the advertised witness:
\begin{equation}
\boxed{
\begin{gathered}
\mathcal F_{\Lambda}=0,
\qquad
F\neq0,\\
\Hol_{A_a}(\ell)\ \text{arbitrary in }\SO(2),
\qquad
\Delta\tau>0.
\end{gathered}}
\label{eq:master}
\end{equation}
The four quantities have different codomains and different physical or geometric duties. Equation \eqref{eq:master} proves coexistence and independent tunability in one model. It supplies a direct witness that exact gluing of local descriptions can close while curvature-bearing physical transport remains nontrivial.

\section{Conclusion}

Compatibility closure is a statement about agreement among local descriptions after the required transition maps are applied. Curvature and holonomy are properties of the physical connection, and proper time is a path functional of the Lorentzian metric. The explicit family above keeps those offices separate by construction. This is the precise sense in which exact gluing can close without imposing flatness.

\begin{acknowledgments}
The author used generative language-model tools for editorial consolidation and algebraic cross-checking. Every definition, computation, boundary, and conclusion was reviewed by the author, who accepts full responsibility for the manuscript.
\end{acknowledgments}

\begin{thebibliography}{9}

\bibitem{KobayashiNomizu1963}
S. Kobayashi and K. Nomizu,
\emph{Foundations of Differential Geometry}, Vol. I
(Wiley, New York, 1963).

\bibitem{Nakahara2003}
M. Nakahara,
\emph{Geometry, Topology and Physics}, 2nd ed.
(Institute of Physics Publishing, Bristol, 2003).

\bibitem{Wald1984}
R. M. Wald,
\emph{General Relativity}
(University of Chicago Press, Chicago, 1984).

\end{thebibliography}

\end{document}
